% Generated draft: requires mathematical and source review.
\documentclass[11pt]{article}
\usepackage[margin=25mm]{geometry}
\usepackage{fontspec}
\setmainfont{Latin Modern Roman}
\usepackage{amsmath,amssymb,mathtools,unicode-math}
\setmathfont{Latin Modern Math}
\usepackage{enumitem,needspace,xurl,hyperref}
\setlength{\parindent}{0pt}
\setlength{\parskip}{5pt}
\setlength{\emergencystretch}{3em}
\newcommand{\N}{\mathbb{N}}
\newcommand{\Z}{\mathbb{Z}}
\newcommand{\Q}{\mathbb{Q}}
\newcommand{\R}{\mathbb{R}}
\newcommand{\C}{\mathbb{C}}
\newcommand{\F}{\mathbb{F}}
\newcommand{\A}{\mathbb{A}}
\newcommand{\PP}{\mathbb{P}}
\newcommand{\E}{\mathbb{E}}
\newcommand{\eps}{\varepsilon}
\newcommand{\dd}{\,\mathrm{d}}
\begin{document}
\section*{UnsolvedMath ID: 30004151}
\textbf{Source problem code:} catalogue entry
\subsection{Definitions and mathematical statement}
A directed acyclic graph is a directed graph with no directed cycle. In the requested causal formulation, an augmented directed acyclic graph is to represent a separation of treatment paths relevant to mediation. The source does not specify the original graph, its augmentation, or the interventions represented by that augmentation. A longitudinal mediator process is a collection of mediator variables indexed by observation times; the time index and the variables observed at each time are unspecified here. The survival outcome concerns time to an event, but the event, observation horizon, and censoring mechanism are unspecified.

For clarity, identifiability can be expressed abstractly. Let \(\mathcal M\) be a specified class of causal models, let \(P_m^{\mathrm{obs}}\) denote the observed-data law induced by \(m\in\mathcal M\), and let \(\theta(m)\) be a specified causal effect. The effect is identifiable over \(\mathcal M\) if
\[
\forall m_1,m_2\in\mathcal M,\qquad
P_{m_1}^{\mathrm{obs}}=P_{m_2}^{\mathrm{obs}}
\ \Longrightarrow\ 
\theta(m_1)=\theta(m_2).
\]
This definition does not supply the missing causal model or effect.

\textbf{Requested research task.} Develop identifiable causal mediation effects for a longitudinal mediator process and a survival outcome by separating treatment paths through an augmented directed acyclic graph. Specify the causal effects and the assumptions under which they are identifiable from the observed-data law. The formulation should address the source's concerns about ill-defined nested counterfactuals and post-treatment confounding. Here, post-treatment confounding refers to confounding involving variables affected by treatment; the source does not identify those variables or their causal relationships.

\textbf{Unresolved formulation.} The supplied record does not determine a unique mathematical proposition. It supplies neither a treatment domain or contrast, a longitudinal observation scheme, a causal graph, a path separation, intervention semantics, nor a survival-effect functional. It also does not give assumptions governing confounding, treatment assignment, mediator assignment, censoring, or observation after the event. Consequently, \(\mathcal M\), \(P_m^{\mathrm{obs}}\), and \(\theta\) above remain placeholders for specifications required by the research task. No particular nested counterfactual is supplied, so the concern about ill-defined counterfactuals cannot be made more specific from this record alone.
\subsection{Short English statement}
Develop causal mediation effects for a mediator evolving over time and a time-to-event outcome, using an augmented causal graph to separate treatment paths. Specify effects and assumptions that permit identification while addressing nested-counterfactual and post-treatment-confounding difficulties. The supplied record leaves the graph, interventions, observation scheme, and effect measure unresolved, so this is a research-task formulation requiring human mathematical review.
\subsection{Sources}
\begin{itemize}
\item Supplied local UnsolvedMath registry record, permanent ID \(30004151\), title: \emph{Identifiable Longitudinal Causal Mediation}; source problem code \texttt{catalogue entry}, in the Oberwolfach Reports Open Problems collection. The record lists proposed year 2019 and supplies no named proposer.
\item Supplied report citation: \emph{Foundations and new horizons for causal inference} (2020), Oberwolfach Reports, DOI 10.4171/owr/2019/25, pp.~1499--1571. \url{https://doi.org/10.4171/owr/2019/25}. The supplied background describes these reports as documenting mathematical workshops and their open-problem discussions. No problem-specific page or passage is supplied.
\item Supplied literature reference, labeled ``Robins--Richardson'': \url{https://arxiv.org/abs/2008.06019}.
\item Supplied literature reference, labeled ``Separable effects'': \url{https://arxiv.org/abs/1901.09472}.
\item Literature assessment dated exactly \texttt{2026-08-22}: the supplied curation labels the problem partially solved and states that interventionist and separable-effects frameworks formalize alternative treatment paths and provide identification results under explicit assumptions. It further states that these frameworks address important mediation settings without resolving the general path-specific causal question. This is supplied dated curation derived from the OWR material, not an independently verified literature assessment; its status and cited claims require independent verification.
\item The registry labels the statement \texttt{corrected\_verified} and reports a manual semantic, typographic, and LaTeX review asserting self-containment and preservation of the mathematical target. These are supplied registry assertions, not verification performed for this draft.
\end{itemize}
\end{document}
